\subsection{Value functions inherit curvature and growth}\label{sec:valuefn}
For the retained set $\mathcal K$ and the recourse set
$\mathcal R=[n]\setminus\mathcal K$ define
\begin{equation}\label{eq:valuefn}
V(v)=\min_{y\in X_{\mathcal R}}F(v,y),\qquad v\in\bar X_{\mathcal K},
\end{equation}
where $\bar X_{\mathcal K}=\prod_{i\in\mathcal K}[\ell_i,u_i]$ and the
feasible set $X_{\mathcal R}$ of the recourse problem is compact and does
not depend on $v$; this is the only structural requirement. For
model~\eqref{eq:model}, $X_{\mathcal R}=\prod_{i\in\mathcal R}X_i$; in
Section~\ref{sec:convexrecourse} it is a product of polytopes. We have
$\min_{X_{\mathcal K}}V=\OPT$, and a retained point $v\in X_{\mathcal K}$
together with a minimizer $y(v)$ of the recourse problem is a feasible point
with value $V(v)$.

\begin{lemma}[Partial minimization]\label{lem:valuefunction}
Let $V$ be defined by~\eqref{eq:valuefn}.
\begin{enumerate}[label=(\alph*)]
\item Let $i\in\mathcal K$ and $L_i\ge0$. If
$t\mapsto F(v+te_i,y)-\frac{L_i}{2}t^2$ is concave on
$\{t:v+te_i\in\bar X_{\mathcal K}\}$ for every $v\in\bar X_{\mathcal K}$
and every $y\in X_{\mathcal R}$, then so is
$t\mapsto V(v+te_i)-\frac{L_i}{2}t^2$ for every $v\in\bar X_{\mathcal K}$.
In particular, upper coordinate curvatures of $F$ for the coordinates in
$\mathcal K$ are upper coordinate curvatures of $V$ on $\bar X_{\mathcal K}$.
\item If $F$ has point growth with constant $g$ at $(v^*,y^*)$, then $V$
has point growth with constant $g$ at $v^*$ on $X_{\mathcal K}$.
\end{enumerate}
\end{lemma}

\begin{proof}
(a) For fixed $v$, the function $t\mapsto V(v+te_i)-\frac{L_i}{2}t^2$ is the
pointwise minimum over $y\in X_{\mathcal R}$ of the concave functions
$t\mapsto F(v+te_i,y)-\frac{L_i}{2}t^2$. It is finite, and it is concave
because its hypograph is the intersection of their hypographs.
(b) Since $\OPT\le V(v^*)\le F(v^*,y^*)=\OPT$, we have $V(v^*)=\OPT$. For
$v\in X_{\mathcal K}$ let $y(v)$ attain the minimum. Then
$V(v)-V(v^*)=F(v,y(v))-\OPT\ge g(\norm{v-v^*}^2+\norm{y(v)-y^*}^2)\ge
g\norm{v-v^*}^2$.
\end{proof}

Part~(a) is the classical stability of semiconcavity under infima
\cite{CannarsaSinestrari2004}. Lemma~\ref{lem:valuefunction}(a) explains why
the corrected-grid bound survives partial minimization, and
Lemma~\ref{lem:valuefunction}(b) why the growth analysis does. The curvature
of $V$ can be much smaller than that of $F$. If $F(v,y)=M(y-v)^2$ with
$v,y\in[0,1]$, the diagonal curvature in $v$ is $2M$, whereas $V\equiv0$.
Section~\ref{sec:convexrecourse} evaluates $V$ exactly and bounds its
curvature by quantities that do not see such stiff convex terms;
Section~\ref{sec:cuts} evaluates $V$ exactly when the eliminated problem is
combinatorial.

The factor structure survives the elimination as follows. Suppose the
recourse coordinates form disjoint \emph{private blocks}
$y^{(1)},\dots,y^{(r)}$ with domains $Y_1,\dots,Y_r$, so that
$X_{\mathcal R}=\prod_tY_t$, and
\begin{equation}\label{eq:privateblocks}
F(v,y)=F_0(v)+\sum_{t=1}^r\phi_t\bigl(v_{E_t},y^{(t)}\bigr),
\end{equation}
where $F_0$ is a sum of factors in $v$ and each block interacts with the
retained coordinates only through its \emph{scope} $E_t\subseteq\mathcal K$.
Then
\[
V(v)=F_0(v)+\sum_{t=1}^r\psi_t(v_{E_t}),\qquad
\psi_t(w)=\min_{y^{(t)}\in Y_t}\phi_t(w,y^{(t)}),
\]
is again a sum of factors, the \emph{value factors} $\psi_t$, with scopes
$E_t$. Any tree decomposition of the scopes of $F_0$ and the sets $E_t$
serves for $V$. Its width is the width after elimination of the blocks,
which can differ from the width of the original decomposition.

\begin{theorem}[Grids over value functions]\label{thm:valuefn}
Suppose $V$ has the form above, with a tree decomposition of the scopes of
$F_0$ and the sets $E_t$ of maximum bag size $p$, and with rational upper
coordinate curvatures $L_i^V$ $(i\in\mathcal K)$ of $V$ whose encoding
lengths are bounded by a polynomial in $I$. Suppose that, for every rational
point $w$, the factors of $F_0$ can be evaluated exactly, and an oracle
returns $\psi_t(w)$ and a minimizer $y^{(t)}(w)$ exactly for every block $t$,
both in time bounded by a polynomial in $I$ and the encoding length of $w$.
Put $L^V=\max_iL_i^V$.
\begin{enumerate}[label=(\alph*)]
\item Theorem~\ref{thm:certificate} holds for $V$, and the incumbent of CT
applied to $V$ lifts to a feasible point $(\hat v,y(\hat v))$ of value
$V(\hat v)$.
\item If $V$ has point growth with constant $g$ (by
Lemma~\ref{lem:valuefunction}(b), this holds if $F$ has), then
Theorem~\ref{thm:approx} holds for $V$ with $\kappa_V=\kappa(L^V,g)$ in
place of $\bar\kappa$ and with the bound
$f(p,\kappa_V)\,\kappa_V^{C}(I+q+1)^{C}$ in place of
$f(p,\bar\kappa)(I+q+1)^5$, where $C$ depends only on the polynomials in
the hypotheses.
\item If $F$ is a rational quadratic on the mixed box $X$
of~\eqref{eq:model}, Theorem~\ref{thm:exact} holds for EX
(Algorithm~\ref{alg:ex}) with CT applied to $V$ and with REC
(Algorithm~\ref{alg:rec}) and the acceptance test applied to $F$ at the
lifted incumbent. If $F$ has point
growth with constant $g$, the bound becomes
$f(p,\kappa_V)\,\kappa_V^{C_1}(I+1)^{C_1}$, where $C_1$ depends only on the
polynomials in the hypotheses.
\end{enumerate}
\end{theorem}

\begin{proof}
(a) Sections~\ref{sec:grids} and~\ref{sec:growth} use the objective only
through Definition~\ref{def:curvature}, exact factor values at grid points,
and, for the running time, growth. The $L_i^V$ satisfy
Definition~\ref{def:curvature} for $V$, the oracle gives exact values and
$y(\hat v)$, and $\min_{X_{\mathcal K}}V=\OPT$.

(b), (c) The accounting of bit lengths, which replaces the common
denominators of Theorem~\ref{thm:approx} by the bounds of the oracle, and of
the accuracy needed for exact output is in
Appendix~\ref{app:recourse-convex}.
\end{proof}

The theorem charges the work of the recourse oracle to the polynomial
factor, so it is useful when the recourse problems are easy and their
elimination does not increase the width. Sections~\ref{sec:convexrecourse}
and~\ref{sec:cuts} give classes in which the recourse problems are solved
exactly and with certificates.
